\section{Body}

\begin{itemize}
\item
	Overview
\item
	Autos
\item
	Security
\item
	Portability
\item
	C comparison
\item
	Perl comparison
\item
	Implementation:
	\begin{itemize}
	\item
		Translation into assembler
	\item
		Compiler optimizations (substituting primitives for autos when possible)
	\end{itemize}

\item
	Penalty reduction:
	\begin{itemize}
	\item
		Convert auto to integer for loop counters.
	\item
		Sometimes alternative primitive code will not be substantially faster anyway.
		For example, if the alternative code is complex and implements a hash with
		reference counts to its elements, clearly it is better to use autos.  This
		is a discussion that people have already written about in terms of Perl vs C,
		so it will have many references.
	\item
		Copy-on-write.  Autos are not copied during assignment or when they are passed by-value to
		functions.  Rather, an alias is created, and copying is delayed until the auto is modified.
		Often the variable will never be modified.  In many cases, this may make autos faster than
		primitive variable types.
	\end{itemize}
\end{itemize}


C Comparison

The following is a C program that adds pairs of numbers read from standard
input.  Following it is the same program written in Pit.  Notice that in
the second, autos manage all the array sizes and memory allocation for us,
thus eliminating any potential for mismanaged string buffers.  As shown in
the third example, they can be managed manually just as easily as in the C
program, with the same risk to security.  This is the trade-off between
improving efficiency and improving security and maintainability.

\sourceinput{comparison.c}

Figure 1.  The program in C.  While this example has no buffer
overflow vulnerabilities, there are several ways they could be
introduced.  For example, the wrong size could be passed to fgets(),
an error could be introduced where the string is parsed with strsep(),
or an error could be introduced in the format string passed to printf.
Also notice that we had to set an arbitrary limit on the length of an
input line.  While it is obviously possible to work around it, doing
so makes the program much more complex and adds more opportunities for
security holes.

\sourceinput{comparison.p}

Figure 2.  The same program, this time written in Pit using autos.
Notice that neither the length of a string nor a pointer into the
string are ever used.

\sourceinput{comparison2.p}

Figure 3.  The same program, this time written in Pit using some
primitive variables.  While autos weren't completely eliminated here,
but they could be.  In particular, \$in is implicitly converted to an
auto when passed to \$string.split, and \$nums remains an auto.  Notice
how easily primitives and autos can work together, allowing the
programmer to choose more precisely what code to optimize for
efficiency at the risk of security, and what code to optimize for
maintainability at the cost of efficiency.
